\chapter{Dunia Kuantum: Foton dan Tingkat Energi}\label{ch:g12:quantum-world}

Sebuah lampu pemanas \omterm{def:g10:light-spectra:wavelength}{inframerah} dapat memandikan kulitmu sepanjang petang tanpa pernah
menandainya; sepuluh menit cahaya matahari gunung yang tipis meninggalkan luka bakar. Lampunya
menyerahkan jauh lebih banyak \omterm{def:g10:energy-conservation:energy}{energi} --- tetapi bukan \omterm{def:g10:energy-conservation:energy}{energi} yang menandai kulit. Cahaya,
akan ditunjukkan bab ini, tiba dalam butiran, dan hanya butiran yang menghantam cukup keras
yang mengerjakan kimia: satu butir ultraungu memutuskan ikatan yang hanya dihangatkan semiliar
butir \omterm{def:g10:light-spectra:wavelength}{inframerah}. Ikutilah butirannya dan sebuah abad terbuka: logam yang membocorkan elektron,
\omterm{def:g11:fundamental-interactions:ladder}{atom} yang bernyanyi dalam \omterm{def:g10:light-spectra:emission}{garis spektrum}, elektron yang \omterm{prop:g12:light-as-wave:conditions}{berinterferensi} bagaikan \omterm{def:g10:signals-and-waves:wave}{gelombang}.

\section{Cahaya tiba dalam butiran: foton}

\begin{definition}[Foton]\label{def:g12:quantum-world:photon}
Cahaya berfrekuensi $\nu$ bertukar \omterm{def:g10:energy-conservation:energy}{energi} dengan benda hanya dalam butiran utuh
yang disebut \emph{foton}\index{foton}, masing-masing berenergi
\[
E = h\nu = \frac{hc}{\lambda}, \qquad h = \qty{6.63e-34}{J.s},
\]
dengan $h$ adalah \emph{tetapan Planck}\index{tetapan Planck} (Planck,
1900; Einstein, 1905) dan $\lambda$ \omterm{def:g12:light-as-wave:lightwave}{panjang gelombangnya di ruang hampa}. Intensitas
sebuah berkas mengukur \emph{banyaknya} foton tiap sekon; \omterm{def:g10:energy-conservation:energy}{energi} tiap butirnya
ditetapkan oleh \omterm{def:g12:oscillators-and-time:oscillator}{frekuensinya} saja.
\end{definition}

\begin{notation}[Elektronvolt]\label{not:g12:quantum-world:ev}
Fisika \omterm{def:g11:fundamental-interactions:ladder}{atom} memakai \emph{elektronvolt}\index{elektronvolt},
$\qty{1}{eV} = \qty{1.60e-19}{J}$: yaitu \omterm{def:g10:energy-conservation:energy}{energi} yang diperoleh satu \omterm{def:g11:fundamental-interactions:elementary}{muatan dasar}
saat melintasi \qty{1}{V} (\cref{ch:g11:electric-gravitational-fields}). Berguna:
$E \,(\unit{eV}) \approx 1240/\lambda\,(\unit{nm})$.
\end{notation}

\begin{omfigure}
\begin{tikzpicture}[font=\small, x=1.48cm]
  \draw[-Stealth] (7.55,0) -- (-0.55,0);
  \draw[-Stealth, omDef] (4.9,-0.85) -- (2.3,-0.85)
    node[midway, below] {energi tiap foton bertambah};
  \foreach \x/\l/\e/\n in {0/{\qty{1}{pm}}/{\qty{1.2}{MeV}}/{$\gamma$},
      1.2/{\qty{0.1}{nm}}/{\qty{12}{keV}}/{X},
      2.4/{\qty{100}{nm}}/{\qty{12}{eV}}/{UV},
      3.6/{\qty{550}{nm}}/{\qty{2.3}{eV}}/{tampak},
      4.8/{\qty{10}{\micro m}}/{\qty{0.12}{eV}}/{IR},
      6.0/{\qty{1}{cm}}/{\qty{0.12}{meV}}/{mikro},
      7.2/{\qty{3}{m}}/{\qty{0.4}{\micro eV}}/{radio}}{
    \draw (\x,0.07) -- (\x,-0.07);
    \node[above, omExo] at (\x,0.42) {\n};
    \node[above, font=\scriptsize] at (\x,0.10) {\l};
    \node[below, font=\scriptsize, omDef] at (\x,-0.10) {\e};}
\end{tikzpicture}

{\small Satu skala bagi keluarga elektromagnetiknya
(\cref{ch:g12:light-as-wave}): \omterm{def:g12:mechanical-waves:wavelength}{panjang gelombangnya} (atas), \omterm{def:g10:energy-conservation:energy}{energi} tiap \omterm{def:g12:quantum-world:photon}{foton}
$hc/\lambda$ (biru), dari butiran radio semikroelektronvolt sampai
butiran \omterm{def:g11:nucleus-radioactivity:abg}{gama} \unit{MeV} pada \cref{ch:g11:nucleus-radioactivity} --- namun sebuah
penunjuk \qty{1.0}{mW} mencurahkan $\num{3e15}$ butir tiap sekon: cahaya
sehari-hari mengalir semulus pasir yang dituang.}
\end{omfigure}

\section{Efek fotolistrik}

Sorotkanlah cahaya pada pelat logam yang bersih di ruang hampa dan elektronnya meletup keluar --- yaitu
\emph{efek fotolistrik}\index{efek fotolistrik}. Kaidahnya mematahkan fisika klasik.

\begin{proposition}[Apa yang ditunjukkan percobaannya]\label{prop:g12:quantum-world:facts}
Bagi tiap logam ada sebuah \emph{frekuensi ambang}\index{frekuensi ambang} $\nu_0$:
\begin{enumerate}
  \item di bawah $\nu_0$, tidak ada elektron yang pergi, seberapa pun kuat cahayanya;
  \item di atas $\nu_0$, pancarannya bermula tanpa tundaan, seberapa pun lemahnya;
  \item intensitasnya menaikkan \emph{banyaknya} elektron tiap sekon, bukan
        \omterm{def:g11:mechanical-energy:kinetic}{energi kinetik} terbesarnya; \omterm{def:g12:oscillators-and-time:oscillator}{frekuensinya} yang menaikkan $E_{k,\max}$.
\end{enumerate}
Tidak satu pun cocok dengan \omterm{def:g10:signals-and-waves:wave}{gelombang} yang malar, yang warna apa pun seharusnya dapat mengisinya ---
selama berbulan-bulan di bawah cahaya yang lemah (\cref{exo:g12:quantum-world:14}) --- dan yang
intensitasnya seharusnya menetapkan \omterm{def:g10:energy-conservation:energy}{energi} elektronnya: ketiga ramalannya gugur.
\end{proposition}
\admitted

\begin{theorem}[Hubungan fotolistrik Einstein]\label{thm:g12:quantum-world:einstein}
\index{hubungan fotolistrik}Sebuah elektron menuntut sekurang-kurangnya \emph{fungsi
kerja}\index{fungsi kerja} $W_0$ logamnya untuk lolos dari permukaannya.
Satu \omterm{def:g12:quantum-world:photon}{foton} diserap utuh oleh satu elektron, sehingga
\[
h\nu = W_0 + E_{k,\max}, \qquad \text{sehingga } \nu_0 = \frac{W_0}{h}.
\]
\end{theorem}

\begin{proof}
Pembukuan \omterm{def:g10:energy-conservation:energy}{energinya}: $h\nu$ masuk, $W_0$ dibelanjakan untuk lolosnya, sisanya kinetik ---
paling banyak $h\nu - W_0$, lebih kecil bagi elektron yang bermula lebih dalam. Di bawah $\nu_0$
satu butir tidak dapat membayar pungutan keluarnya, berapa pun banyaknya.
\end{proof}

\begin{definition}[Potensial henti]\label{def:g12:quantum-world:stopping}
Untuk mengukur $E_{k,\max}$, \omterm{def:g11:circuits-and-power:voltage}{tegangan} balik mengerem elektronnya
(\cref{ch:g11:electric-gravitational-fields}); \emph{potensial
henti}\index{potensial henti} $U_s$ adalah \omterm{def:g11:circuits-and-power:voltage}{tegangan} terkecil yang meniadakan
arusnya: elektron yang tercepat tepat gagal menyeberang, sehingga $E_{k,\max} = e\,U_s$.
\end{definition}

\begin{omfigure}
\begin{tikzpicture}
\begin{axis}[omaxis, width=8.8cm, height=5.4cm,
    xlabel={$\nu$ ($10^{14}\,\unit{Hz}$)}, ylabel={$E_{k,\max}$ (\unit{eV})},
    xmin=0, xmax=13.4, ymin=-3.1, ymax=3.2,
    xtick={2,4,6,8,10,12}, ytick={-2,-1,1,2,3}]
  \addplot[omExo, dashed, domain=0:5.55, samples=2] {0.414*x-2.3};
  \addplot[omThm, thick, domain=5.55:12.4, samples=2] {0.414*x-2.3};
  \addplot[only marks, mark=*, mark size=1.5pt, omDef] coordinates
    {(7,0.60) (8,1.01) (9,1.43) (10,1.84) (11,2.26)};
  \node[below, font=\small] at (axis cs:5.55,-0.08) {$\nu_0$};
  \node[right, font=\small, omExo] at (axis cs:0.15,-2.55) {$-W_0$};
  \node[font=\small, omThm, rotate=22] at (axis cs:10.3,1.35) {kemiringan $= h$};
\end{axis}
\end{tikzpicture}

{\small $E_{k,\max} = eU_s$ terukur terhadap \omterm{def:g12:oscillators-and-time:oscillator}{frekuensinya} (natrium): sebuah
garis lurus berkemiringan $h$, yang memotong sumbunya di $\nu_0$; jika diperpanjang
(bergaris putus-putus), ia menyingkapkan $-W_0$ pada $\nu = 0$.}
\end{omfigure}

\begin{method}[Mengukur tetapan Planck]\label{met:g12:quantum-world:measure-h}
\begin{enumerate}
  \item Sinarilah sebuah fotosel melalui tapis yang \omterm{def:g12:mechanical-waves:wavelength}{panjang gelombangnya} diketahui;
        hitunglah tiap $\nu = c/\lambda$; catatlah tiap \omterm{def:g12:quantum-world:stopping}{potensial hentinya}.
  \item Lukiskanlah $E_{k,\max} = eU_s$ terhadap $\nu$: Einstein menjanjikan
        garis lurus $h\nu - W_0$.
  \item Bacalah kemiringannya: itu \emph{adalah} $h$; perpotongannya memberikan $\nu_0$ dan $-W_0$.
\end{enumerate}
\end{method}

\section{Tingkat energi: mengapa atom memancarkan garis}

\begin{definition}[Tingkat energi]\label{def:g12:quantum-world:levels}
\omterm{def:g10:energy-conservation:energy}{Energi} sebuah elektron yang terikat di dalam \omterm{def:g11:fundamental-interactions:ladder}{atom} bersifat
\emph{terkuantum}\index{kuantisasi}: terbatas pada tangga diskret berisi
\emph{tingkat energi}\index{tingkat energi} $E_1 < E_2 < \dots < 0$, yang dihitung
negatif dari elektron bebasnya. Anak tangga terendahnya adalah \emph{keadaan
dasar}\index{keadaan dasar}, yang lain \emph{keadaan
tereksitasi}\index{keadaan tereksitasi}; mencapai $0$ adalah
\emph{ionisasi}\index{ionisasi}.
\end{definition}

\begin{proposition}[Tangga hidrogennya]\label{prop:g12:quantum-world:hydrogen}
\omterm{def:g12:quantum-world:levels}{Tingkat energi} \omterm{def:g11:fundamental-interactions:ladder}{atom} hidrogen adalah
\[
E_n = -\frac{\qty{13.6}{eV}}{n^2}, \qquad n = 1, 2, 3, \dots
\]
sehingga $E_1 = \qty{-13.6}{eV}$, $E_2 = \qty{-3.40}{eV}$, $E_3 =
\qty{-1.51}{eV}$, $E_4 = \qty{-0.85}{eV}$, yang berdesakan menuju $0$.
\end{proposition}
\admitted

\begin{remark}[Dari mana tangganya datang]\label{rem:g12:quantum-world:ladder-origin}
Mengapa bilangan inilah yang muncul adalah mekanika kuantum, yang disimpan untuk jilid universitasnya ---
begitu pula rumus de Broglie: \omterm{def:g12:quantum-world:debroglie}{gelombang materi} elektronnya
(\cref{def:g12:quantum-world:debroglie} di bawah) harus menutup pada dirinya sendiri mengelilingi
intinya, dan hanya \omterm{def:g12:mechanical-waves:wavelength}{panjang gelombang} tertentu yang muat --- bagaikan senar gitar, sebuah
\omterm{def:g11:fundamental-interactions:ladder}{atom} memiliki himpunan nada yang diskret.
\end{remark}

\begin{theorem}[Foton sebuah transisi]\label{thm:g12:quantum-world:transition}
Sebuah \omterm{def:g11:fundamental-interactions:ladder}{atom} berpindah tingkat hanya dengan memancarkan atau menyerap satu \omterm{def:g12:quantum-world:photon}{foton} yang membawa
tepat selisihnya: turun dari tingkat $p$ ke tingkat $n$ yang lebih rendah memancarkan
$h\nu = E_p - E_n$; menanjak dari $n$ ke $p$ menyerap \omterm{def:g10:energy-conservation:energy}{energi} yang sama itu.
\end{theorem}

\begin{proof}
\omterm{thm:g10:energy-conservation:conservation}{Kekekalan energi}, satu butir pada satu waktu: fotonlah satu-satunya mata uang
yang diterima cahaya, dan pembukuan \omterm{def:g11:fundamental-interactions:ladder}{atomnya} harus setimbang.
\end{proof}

\begin{remark}[Spektrum garis yang terjelaskan]\label{rem:g12:quantum-world:lines}
Inilah jawaban yang dijanjikan \cref{ch:g10:light-spectra}: gas yang panas memancarkan
hanya \omterm{def:g12:oscillators-and-time:oscillator}{frekuensi} $(E_p - E_n)/h$ milik tangganya sendiri --- sebuah kode batang
berisi garis terang --- dan gas yang dingin mencuri tepat \omterm{def:g12:oscillators-and-time:oscillator}{frekuensi} itu dari \omterm{def:g10:light-spectra:white}{cahaya
putih}, mencetak garis serapan gelap yang bersesuaian. Tiap unsur memiliki tangganya
sendiri, sehingga kode batangnya sendiri --- begitulah kita membaca cahaya bintang.
\end{remark}

\begin{example}[Garis tampak hidrogen]\label{ex:g12:quantum-world:balmer}
Transisi yang berakhir pada $n = 2$ jatuh di daerah tampak. Dengan $E \,(\unit{eV}) =
1240/\lambda\,(\unit{nm})$: $3 \to 2$: \qty{1.89}{eV}, \qty{656}{nm} (merah);
$4 \to 2$: \qty{2.55}{eV}, \qty{486}{nm} (biru kehijauan); $5 \to 2$:
\qty{2.86}{eV}, \qty{434}{nm} (ungu). Transisi ke $n = 1$ melampaui
\qty{10}{eV}: semuanya ultraungu.
\end{example}

\begin{omfigure}
\begin{tikzpicture}[font=\small, y=1cm]
  \foreach \y/\n/\e in {0/1/-13.6, 4.50/2/-3.40, 5.33/3/-1.51}{
    \draw[thick] (0,\y) -- (4.2,\y);
    \node[left] at (0,\y) {$n=\n$};
    \node[right, font=\scriptsize] at (4.2,\y) {\qty{\e}{eV}};}
  \draw[thick] (0,5.63) -- (4.2,5.63) node[left,at start,yshift=-2.5pt] {$n=4$};
  \draw[thick] (0,5.76) -- (4.2,5.76) node[left,at start,yshift=3.5pt] {$n=5$};
  \draw[dashed] (0,6.0) -- (4.2,6.0) node[left, at start] {$\infty$}
    node[right, font=\scriptsize] {\qty{0}{eV} (terionkan)};
  \foreach \x/\ys/\col/\l in {1.0/5.33/omThm/656, 1.9/5.63/omDef/486,
      2.8/5.76/violet/434} \draw[-Stealth, \col, thick] (\x,\ys) --
    (\x,4.50) node[midway, left, font=\scriptsize] {\qty{\l}{nm}};
  \draw[-Stealth, omExo, thick] (3.7,4.50) -- (3.7,0)
    node[midway, right, font=\scriptsize] {\qty{122}{nm} (UV)};
\end{tikzpicture}

{\small Tangga hidrogennya, sesuai skala: garis tampaknya semua berakhir pada
$n = 2$; turunnya ke \omterm{def:g12:quantum-world:levels}{keadaan dasar} berada jauh di ultraungu.}
\end{omfigure}

\begin{remark}[Laser, dalam satu tarikan napas]\label{rem:g12:quantum-world:laser}
Sebuah \omterm{def:g12:quantum-world:photon}{foton} berenergi tepat $E_p - E_n$ yang melewati \omterm{def:g11:fundamental-interactions:ladder}{atom} yang \emph{tereksitasi} dapat memicunya
memancarkan \omterm{def:g12:quantum-world:photon}{foton} yang serupa --- \omterm{def:g10:energy-conservation:energy}{energi} yang sama, arah yang sama, selaras. Jagalah
lebih banyak \omterm{def:g11:fundamental-interactions:ladder}{atomnya} di atas daripada di bawah dan satu \omterm{def:g12:quantum-world:photon}{foton} menjadi longsoran klon:
\emph{pancaran terangsang}\index{pancaran terangsang} --- dan \omterm{def:g11:color-light-sources:lamps}{laser} tidak lebih dari itu.
\end{remark}

\section{Gelombang materi}

\begin{proposition}[Elektron berinterferensi]\label{prop:g12:quantum-world:interference}
Kirimkanlah elektron satu per satu melalui celah ganda
(\cref{ch:g12:light-as-wave}): masing-masing mendarat sebagai satu titik yang tajam --- sebuah
zarah. Tetapi sembari ribuan menumpuk, titiknya melukis rumbai interferensi:
tiap elektron, sendirian, melintas sebagai \omterm{def:g10:signals-and-waves:wave}{gelombang} melalui \emph{kedua} celahnya. Hablur
melenturkan berkas elektron juga (Davisson dan Germer, 1927).
\end{proposition}
\admitted

\begin{omfigure}
\begin{tikzpicture}[font=\small]
  \foreach \o in {0,4.0,8.0} \draw[omExo] (\o,0) rectangle +(3.4,2.2);
  \foreach \x/\y in {0.62/1.71, 1.68/0.42, 2.31/1.95, 1.05/1.22,
      2.98/0.66, 1.74/1.58, 0.38/0.35, 2.40/0.98, 1.62/1.02, 0.97/1.86,
      2.95/1.60, 1.71/0.20, 1.12/0.61, 2.33/0.31}
    \fill (\x,\y) circle (0.8pt);
  \foreach \c/\f in {0.40/610, 1.05/410, 1.70/530, 2.35/350, 3.00/470}
    \foreach \y in {0.25,0.62,...,2.00}
      \fill ({4.0+\c+0.14*sin(\f*\y)},\y) circle (0.8pt);
  \foreach \x/\y in {4.72/1.42, 5.38/0.52, 6.06/1.12, 6.66/1.72,
      5.44/1.88, 7.28/0.78} \fill (\x,\y) circle (0.8pt);
  \foreach \c/\f in {0.40/610, 1.05/410, 1.70/530, 2.35/350, 3.00/470}
    \foreach \y in {0.12,0.27,...,2.10}
      \fill ({8.0+\c+0.11*sin(\f*\y)},\y) circle (0.8pt);
  \foreach \c/\f in {0.40/260, 1.05/340, 1.70/205, 2.35/290, 3.00/385}
    \foreach \y in {0.20,0.50,...,2.05}
      \fill ({8.0+\c+0.06*sin(\f*\y+40)},\y) circle (0.8pt);
  \node[below] at (1.7,-0.05) {$\sim 30$ elektron};
  \node[below] at (5.7,-0.05) {$\sim 200$}; \node[below] at (9.7,-0.05) {$\sim 2000$};
\end{tikzpicture}

{\small Satu elektron, satu titik; ribuan, rumbai. Di mana titik berikutnya mendarat
bersifat acak --- seperti peluruhan satu inti (\cref{ch:g11:nucleus-radioactivity})
--- tetapi \omterm{def:g10:signals-and-waves:wave}{gelombangnya} menetapkan \emph{peluangnya}: rumbai terang, besar peluangnya; gelap,
hampir tidak pernah. Kebetulan perorangan yang diatur \omterm{def:g10:signals-and-waves:wave}{gelombang} yang tepat: jantung fisika kuantum.}
\end{omfigure}

\begin{definition}[Panjang gelombang de Broglie]\label{def:g12:quantum-world:debroglie}
Setiap zarah bermomentum $p = mv$ berjalan sebagai \emph{gelombang
materi}\index{gelombang materi} ber\emph{panjang gelombang de Broglie}\index{panjang gelombang de
Broglie}
\[
\lambda = \frac{h}{p} = \frac{h}{mv}.
\]
\end{definition}
\admitted

\begin{example}[Mengapa kamu tidak pernah melentur]\label{ex:g12:quantum-world:never}
Sebuah elektron yang dipercepat melalui \qty{100}{V} memiliki $E_k = \qty{1.60e-17}{J}$,
$p = \sqrt{2m_e E_k} = \qty{5.4e-24}{kg.m/s}$, sehingga $\lambda \approx
\qty{0.12}{nm}$ --- yaitu jarak antaratom di dalam hablur, yang karenanya bertindak
sebagai celah gandanya. Sebuah bola tenis \qty{58}{g} pada \qty{25}{m/s}: $\lambda =
6.63\times10^{-34}/1.45 \approx \qty{4.6e-34}{m}$, sekitar $10^{19}$ kali
lebih kecil daripada sebuah inti --- tidak akan ada celah yang memperlihatkannya berumbai. $h$ begitu
kecil sehingga kebutiran kuantumnya muncul hanya pada skala \omterm{def:g11:fundamental-interactions:ladder}{atomnya}.
\end{example}

\begin{example}[Mikroskop elektron]\label{ex:g12:quantum-world:microscope}
Pelenturan melarang memilah rincian yang jauh lebih kecil daripada \omterm{def:g12:mechanical-waves:wavelength}{panjang gelombang} yang dipakai
(\cref{ch:g12:light-as-wave}): cahaya tampak berhenti di dekat \qty{0.5}{\micro m}.
Elektron pada \qty{10}{kV} memiliki $\lambda \approx \qty{12}{pm}$ --- puluhan
ribu kali lebih pendek: mikroskop elektron mencitrakan virus, mesin
sel, bahkan \omterm{def:g11:fundamental-interactions:ladder}{atom} tunggal. \omterm{def:g12:quantum-world:debroglie}{Gelombang materi} menyerahkan sebuah mata yang baru kepada fisika.
\end{example}

\section{Latihan}

\begin{exercise}[$\star$]\label{exo:g12:quantum-world:1}
Sebuah penunjuk \omterm{def:g11:color-light-sources:lamps}{laser} hijau memancarkan \qty{1.0}{mW} pada \qty{530}{nm}. Hitunglah
\omterm{def:g10:energy-conservation:energy}{energi} satu \omterm{def:g12:quantum-world:photon}{foton} dalam \omterm{def:g11:work-of-force:work}{joule} dan \omterm{def:g12:nuclear-energy:units}{elektronvolt}, lalu banyaknya \omterm{def:g12:quantum-world:photon}{foton} tiap sekon.
\end{exercise}

\begin{exercise}[$\star$]\label{exo:g12:quantum-world:2}
Sebuah stasiun FM memancar pada \qty{100}{MHz}. Hitunglah \omterm{def:g10:energy-conservation:energy}{energi} \omterm{def:g12:quantum-world:photon}{fotonnya} dalam
\unit{eV} lalu bandingkan dengan \omterm{def:g12:quantum-world:photon}{foton} tampak: mengapa kebutiran radionya tak terdeteksi?
\end{exercise}

\begin{exercise}[$\star$]\label{exo:g12:quantum-world:3}
Sesium memiliki $W_0 = \qty{1.9}{eV}$. Hitunglah \omterm{def:g12:oscillators-and-time:oscillator}{frekuensi} dan \omterm{def:g12:mechanical-waves:wavelength}{panjang gelombang}
ambangnya. Apakah \omterm{def:g11:color-light-sources:lamps}{laser} merah pada \qty{633}{nm} melontarkan elektron dari sesium ---
dan jika ya, dengan \omterm{def:g11:mechanical-energy:kinetic}{energi kinetik} terbesar berapa?
\end{exercise}

\begin{exercise}[$\star$]\label{exo:g12:quantum-world:4}
Dengan $E_n = -13.6/n^2$ \unit{eV}, hitunglah \omterm{def:g10:energy-conservation:energy}{energi} dan \omterm{def:g12:mechanical-waves:wavelength}{panjang gelombang}
\omterm{def:g12:quantum-world:photon}{foton} yang dipancarkan pada transisi hidrogen $3 \to 2$. Apa warnanya?
\end{exercise}

\begin{exercise}[$\star$]\label{exo:g12:quantum-world:5}
Dari grafik $E_{k,\max}$ terhadap $\nu$ pada bab ini, bacalah \omterm{prop:g12:quantum-world:facts}{frekuensi ambang}
dan \omterm{thm:g12:quantum-world:einstein}{fungsi kerja} natrium (dalam \unit{eV}); tetapan apa yang menjadi kemiringannya?
\end{exercise}

\begin{exercise}[$\star\star$]\label{exo:g12:quantum-world:6}
Cahaya berpanjang \omterm{def:g10:signals-and-waves:wave}{gelombang} \qty{400}{nm} jatuh pada kalium ($W_0 = \qty{2.3}{eV}$).
Hitunglah \omterm{def:g10:energy-conservation:energy}{energi} \omterm{def:g12:quantum-world:photon}{fotonnya}, \omterm{def:g11:mechanical-energy:kinetic}{energi kinetik} terbesar elektronnya
(dalam \unit{eV} dan \omterm{def:g11:work-of-force:work}{joule}), dan \omterm{def:g12:quantum-world:stopping}{potensial hentinya}.
\end{exercise}

\begin{exercise}[$\star\star$]\label{exo:g12:quantum-world:7}
Sebuah fotosel disinari di atas ambangnya. Apa yang terjadi pada (a) arusnya, (b)
\omterm{def:g12:quantum-world:stopping}{potensial hentinya}, ketika intensitasnya digandakan pada \omterm{def:g12:oscillators-and-time:oscillator}{frekuensi} yang tetap? Ketika
\omterm{def:g12:oscillators-and-time:oscillator}{frekuensinya} dinaikkan? Fakta yang mana yang bertentangan dengan gambaran \omterm{def:g10:signals-and-waves:wave}{gelombangnya}, dan mengapa?
\end{exercise}

\begin{exercise}[$\star\star$]\label{exo:g12:quantum-world:8}
\omterm{def:g10:energy-conservation:energy}{Energi} \omterm{def:g12:quantum-world:photon}{foton} terkecil berapa yang mengionkan \omterm{def:g11:fundamental-interactions:ladder}{atom} hidrogen dari \omterm{def:g12:quantum-world:levels}{keadaan dasarnya}? Dari
$n = 2$? Hitunglah kedua \omterm{def:g12:mechanical-waves:wavelength}{panjang gelombang} ambangnya lalu sebutkan daerah \omterm{def:g12:sound-acoustics:timbre}{spektrumnya}.
\end{exercise}

\begin{exercise}[$\star\star$]\label{exo:g12:quantum-world:9}
\omterm{def:g10:light-spectra:white}{Cahaya putih} melintasi sebuah toples berisi hidrogen yang dingin (semua \omterm{def:g11:fundamental-interactions:ladder}{atomnya} di \omterm{def:g12:quantum-world:levels}{keadaan dasar}).
Jelaskan mengapa \omterm{def:g12:sound-acoustics:timbre}{spektrum} yang diteruskan \emph{tidak} menunjukkan garis gelap di
daerah tampak, namun garis serapan Balmer muncul pada \omterm{def:g12:sound-acoustics:timbre}{spektrum} bintang yang panas.
\end{exercise}

\begin{exercise}[$\star\star$]\label{exo:g12:quantum-world:10}
Hitunglah \omterm{def:g12:quantum-world:debroglie}{panjang gelombang de Broglie} (a) sebuah elektron yang dipercepat melalui
\qty{100}{V}, (b) seekor lalat \qty{0.10}{g} pada \qty{1.0}{m/s}. Bandingkanlah masing-masing dengan
ukuran sebuah \omterm{def:g11:fundamental-interactions:ladder}{atom} ($\sim \qty{e-10}{m}$); lalu simpulkan.
\end{exercise}

\begin{exercise}[$\star\star$]\label{exo:g12:quantum-world:11}
Mata yang terbiasa gelap menanggapi sekitar $90$ \omterm{def:g12:quantum-world:photon}{foton} cahaya \qty{505}{nm}.
Hitunglah \omterm{def:g10:energy-conservation:energy}{energi} itu, bandingkan dengan \omterm{def:g11:mechanical-energy:kinetic}{energi kinetik} seekor nyamuk \qty{2.5}{mg}
yang terbang pada \qty{0.50}{m/s}, lalu tafsirkanlah nisbahnya.
\end{exercise}

\begin{exercise}[$\star\star\star$]\label{exo:g12:quantum-world:12}
Sebuah fotosel memberikan $U_s = \qty{1.19}{V}$ pada \qty{405}{nm} dan $U_s =
\qty{0.40}{V}$ pada \qty{546}{nm}. Dari kedua titik ini ukurlah $h$, lalu
\omterm{thm:g12:quantum-world:einstein}{fungsi kerjanya} dalam \unit{eV}. Logam mana pada bab ini yang menjadi katodanya?
\end{exercise}

\begin{exercise}[$\star\star\star$]\label{exo:g12:quantum-world:13}
Sebuah lampu hidrogen menunjukkan garis pada \qty{486}{nm} dan \qty{434}{nm}. Ubahlah
masing-masing menjadi \omterm{def:g10:energy-conservation:energy}{energi} \omterm{def:g12:quantum-world:photon}{foton} lalu kenalilah kedua transisinya dengan mencocokkan
selisih tingkat $E_n = -13.6/n^2$ \unit{eV}.
\end{exercise}

\begin{exercise}[$\star\star\star$]\label{exo:g12:quantum-world:14}
Taksiran tundaan klasiknya: cahaya lemah berintensitas \qty{1.0e-6}{W/m^2}
jatuh pada kalium ($W_0 = \qty{2.3}{eV}$). Jika sebuah \omterm{def:g11:fundamental-interactions:ladder}{atom} hanya memungut
\omterm{def:g10:energy-conservation:energy}{energi} yang melintasi tampang lintangnya (berjari-jari \qty{1.0e-10}{m}), berapa lama
ia akan menuntut untuk mengumpulkan $W_0$? Bandingkan dengan tundaan teramatinya (di bawah \qty{e-9}{s}).
\end{exercise}

\begin{exercise}[$\star\star\star$]\label{exo:g12:quantum-world:15}
Di dalam mikroskop elektron, elektronnya dipercepat melalui \qty{5.0}{kV}.
Hitunglah \omterm{def:g12:mechanical-waves:wavelength}{panjang gelombang} de Brogliennya, bandingkan dengan cahaya hijau (\qty{550}{nm}),
lalu jelaskan lewat pelenturan (\cref{ch:g12:light-as-wave}) mengapa ini membeli \omterm{def:g11:work-of-force:power}{daya} pilah.
\end{exercise}

\section{Soal: Laboratorium fisika panel surya}

\begin{problem}[{Soal akhir pekan --- laboratorium fisika panel surya:
$10^{21}$ butir yang diminum sebuah atap tiap sekon, sebuah fotosel yang mengukur
\omterm{def:g12:quantum-world:photon}{tetapan Planck}, sebuah lampu hidrogen yang dibaca garis demi garis --- satu akhir pekan, satu $h$}]
\label{pb:g12:quantum-world:1}
Kelasmu memasangi panel surya atap sekolahmu dengan alat selama satu akhir pekan.
Datanya: iradiansi surya pada tengah hari \qty{1.0e3}{W/m^2}; \omterm{def:g12:mechanical-waves:wavelength}{panjang gelombang} surya rata-rata
\qty{550}{nm}; luas panelnya \qty{1.0}{m^2}, dayagunanya $20\%$; harinya
menyerahkan setara \qty{5.0}{h} matahari penuh.

\medskip
\noindent\textbf{Bagian I --- Pembukuan \omterm{def:g12:quantum-world:photon}{fotonnya}.}
\begin{enumerate}
  \item Hitunglah \omterm{def:g10:energy-conservation:energy}{energi} satu \omterm{def:g12:quantum-world:photon}{foton} surya rata-rata, dalam \omterm{def:g11:work-of-force:work}{joule} dan \omterm{def:g12:nuclear-energy:units}{elektronvolt}.
  \item Perolehlah banyaknya \omterm{def:g12:quantum-world:photon}{foton} yang menghantam panelnya tiap sekon pada tengah hari.
  \item Hitunglah \omterm{prop:g11:circuits-and-power:power}{daya listrik} panelnya pada tengah hari, lalu \omterm{def:g10:energy-conservation:energy}{energi} yang
        dihasilkannya sepanjang hari, dalam \unit{kW.h}.
  \item Arus panelnya tampak mulus sempurna. Dengan memakai pertanyaan 2,
        jelaskan mengapa kedatangan butir demi butirnya tak terlihat.
  \item Cahaya matahari membawa lebih banyak \omterm{def:g10:energy-conservation:energy}{energi} \omterm{def:g10:light-spectra:wavelength}{inframerah} daripada ultraungu, namun hanya
        ultraungu yang membakar kulit (ikatannya $\sim \qtyrange{3}{4}{eV}$). Jelaskan dengan butirannya.
\end{enumerate}

\medskip
\noindent\textbf{Bagian II --- Mengukur $h$ dengan fotosel perkakasnya.}
Fotosel perkakasnya berkatoda tak dikenal. Tapisnya memberikan:
\begin{center}\small
\begin{tabular}{lcccc}
$\lambda$ (\unit{nm}) & 365 & 405 & 436 & 546\\
$U_s$ (\unit{V}) & 1.50 & 1.16 & 0.94 & 0.37
\end{tabular}
\end{center}
\begin{enumerate}[resume]
  \item Apa yang diukur \omterm{def:g12:quantum-world:stopping}{potensial hentinya}? Hubungkanlah $U_s$, $\nu$, $W_0$, $h$.
  \item Ubahlah keempat \omterm{def:g12:mechanical-waves:wavelength}{panjang gelombangnya} menjadi \omterm{def:g12:oscillators-and-time:oscillator}{frekuensi} (dalam $10^{14}\,\unit{Hz}$).
  \item Mengapa melukiskan $eU_s$ terhadap $\nu$ seharusnya memberikan garis lurus?
        Periksalah kesegarisan titiknya.
  \item Dari titik-titik ujungnya, hitunglah kemiringannya: itulah $h$ hasil ukuranmu.
  \item Perolehlah \omterm{thm:g12:quantum-world:einstein}{fungsi kerja} katodanya (\unit{eV}), \omterm{prop:g12:quantum-world:facts}{frekuensi
        ambangnya} dan \omterm{def:g12:mechanical-waves:wavelength}{panjang gelombang} ambangnya.
  \item Bagian mana pada cahaya mataharinya yang dapat melontarkan elektron dari
        katoda ini --- dan separuh \omterm{def:g12:sound-acoustics:timbre}{spektrum} yang mana yang tak berguna baginya?
  \item Bandingkanlah $h$-mu dengan \qty{6.63e-34}{J.s}: berapa galat nisbinya dalam persen?
\end{enumerate}

\medskip
\noindent\textbf{Bagian III --- Lampu peneraan hidrogennya.}
Skala \omterm{def:g12:mechanical-waves:wavelength}{panjang gelombang} perkakasnya diperiksa terhadap sebuah lampu hidrogen,
$E_n = -13.6/n^2$ \unit{eV}.
\begin{enumerate}[resume]
  \item Hitunglah $E_2$, $E_3$, $E_4$ dan $E_5$.
  \item Hitunglah \omterm{def:g10:energy-conservation:energy}{energi} \omterm{def:g12:quantum-world:photon}{foton} transisi $3 \to 2$, $4 \to 2$, $5 \to 2$.
  \item Perolehlah ketiga \omterm{def:g12:mechanical-waves:wavelength}{panjang gelombangnya} lalu sebutkan warna tiap garisnya.
  \item Mengapa lampunya memancarkan \emph{garis} dan bukan pelangi yang
        malar? Satu kalimat, dengan menyebutkan gagasan kuncinya.
  \item Dapatkah hidrogen lampunya menyerap cahaya pada \qty{500}{nm}? Berikan alasannya dengan tangganya.
\end{enumerate}

\medskip
\noindent\textbf{Bagian IV --- Memeriksa selnya: \omterm{def:g10:signals-and-waves:wave}{gelombang} elektronnya.}
Retak mikro di dalam sebuah sel terletak jauh di bawah capaian \qty{0.5}{\micro m} mikroskop
cahaya; mikroskop elektron payar laboratoriumnya berjalan pada \qty{10}{kV}.
\begin{enumerate}[resume]
  \item Hitunglah laju elektronnya; periksalah bahwa lajunya tetap di bawah $0.25\,c$ (rumus klasik).
  \item Hitunglah \omterm{def:g12:newtons-laws:momentum}{momentumnya} dan \omterm{def:g12:mechanical-waves:wavelength}{panjang gelombang} de Brogliennya.
  \item \omterm{def:g10:energy-conservation:power}{Daya} pilahnya berskala menurut \omterm{def:g12:mechanical-waves:wavelength}{panjang gelombangnya}: hitunglah perolehannya atas cahaya hijau
        (\qty{550}{nm}); simpulkanlah dengan kedua bilangan akhir pekannya --- $h$-mu
        dan perolehan itu.
\end{enumerate}
\end{problem}
